\chapter{Introduction}
\label{ch:intro}

Automated surface quality assurance in high-throughput food container and packaging manufacturing represents a critical engineering imperative to guarantee structural hermeticity, preserve shelf-life integrity, and safeguard consumer health. In modern food manufacturing facilities, packaging containers---specifically thermoformed and injection-molded polymer casings utilized for ready-to-eat rice meals---are manufactured at rates exceeding several thousands of units per hour. During these high-speed cyclic thermal and mechanical forming operations, polymeric casings are susceptible to severe morphological, physical, and particulate anomalies. These flaws encompass polymer carbonization black dots, guillotine edge-trimming tears, micro-abrasion scratches, sub-millimeter thermoforming pinholes, structural shrivel wrinkles, and biological contaminants such as insects.

Among these anomalies, micro-perforations spanning less than 0.2~mm in diameter and elongated diagonal scratches compromise the protective atmospheric barrier of the packaging. Such breaches allow oxygen ingress, microbial infiltration, and moisture migration, precipitating rapid food spoilage, commercial product recalls, and severe brand liability. Ensuring complete zero-defect escape in high-speed packaging requires an inspection mechanism capable of detecting, classifying, and precisely localizing multi-scale anomalies at line speeds exceeding 1.2~m/s. However, the industrial environment introduces harsh constraints: variable ambient illumination, non-uniform specular reflections across curved polymeric surfaces, high optical conveyance jitter, and strict edge computational latency budgets.

This project formulates, develops, and experimentally validates a vision-based defect detection system for industrial quality inspection driven by an optimized YOLOv8 Oriented Bounding Box (YOLOv8-OBB) architecture to overcome these industrial inspection bottlenecks~\cite{wang2024rgs,terven2023yolo}.

%------------------------------------------------------------------
\section{Background}
\label{sec:background}

The food packaging sector constitutes a multi-billion-dollar global enterprise operating under stringent international sanitary and quality standards (e.g., ISO 22000, HACCP, FDA Title 21). Packaging casings constructed from multi-layer polypropylene (PP), polyethylene terephthalate (PET), and high-barrier ethylene vinyl alcohol (EVOH) copolymers serve as the primary line of defense against biological decomposition~\cite{patel2023deep}. As high-speed automated packaging machinery scales production velocities to meet global consumer demand, the temporal window available for visual inspection shrinks to tens of milliseconds per unit. Consequently, manual visual inspection and legacy rule-based optical inspection systems have become major operational bottlenecks in modern packaging facilities.

\subsection{Existing Methods}
\label{subsec:existing}

Industrial quality control on packaging conveyance lines has traditionally relied upon three foundational inspection paradigms:

\begin{itemize}
\item \textbf{Manual Human Visual Inspection:} Human inspectors observe moving conveyor lines to manually detect and segregate anomalous containers. This approach remains widely deployed across many food and consumer packaging plants despite severe speed and accuracy limitations.
\item \textbf{Traditional Rule-Based Automated Optical Inspection (AOI):} First-generation automated optical inspection systems employ handcrafted image processing operators, including Sobel and Canny edge detectors, Otsu dynamic thresholding, morphological opening/closing filters, and template matching algorithms. These algorithms rely on rigid photometric and spatial assumptions to identify contrast variations on container rims and surfaces.
\item \textbf{Standard Deep Learning Horizontal Detectors (HBB):} In recent years, standard deep learning object detectors such as Faster R-CNN, SSD, YOLOv5, and default YOLOv8 have been introduced to industrial inspection. These detectors employ standard horizontal axis-aligned bounding boxes (HBB) defined by coordinates $[x_{\min}, y_{\min}, x_{\max}, y_{\max}]$ to delineate defective regions~\cite{terven2023yolo,gupta2024benchmark}.
\end{itemize}

\subsection{Limitations of Existing Methods}
\label{subsec:limitations}

Despite their widespread historical adoption, existing inspection methodologies suffer from severe operational limitations:

\begin{itemize}
\item \textbf{Physiological Fatigue and Throughput Ceilings:} Human inspection suffers from rapid physiological eye fatigue, cognitive lapses, subjectivity in defect grading, and an inability to sustain inspection accuracy on conveyor lines moving faster than 0.5~m/s. Empirical industrial studies reveal defect escape rates of 12\% to 18\% during sustained human shifts.
\item \textbf{Environmental Sensitivity and False Alarm Spikes:} Classical rule-based algorithms are notoriously brittle when deployed in dynamic factory environments. Curved thermoformed casings exhibit non-uniform specular reflections, subtle color variations across different polymer batches, and ambient factory lighting fluctuations. Rule-based edge filters frequently trigger false-positive alarms on harmless surface reflections while missing low-contrast scratches and pinholes.
\item \textbf{Severe Feature Dilution in Axis-Aligned Bounding Boxes:} Standard horizontal bounding boxes (HBB) present a fundamental geometric mismatch when applied to slender, arbitrarily oriented industrial flaws such as diagonal scratches, curving flange tears, and oblique shrivel wrinkles. An axis-aligned rectangle enclosing a 45-degree scratch inevitably encloses up to 85\% extraneous non-defective background pixels. This causes severe feature dilution, where the neural network's convolutional layers average defect gradients with background noise, leading to catastrophic miss rates (empirically exceeding 87\% for fine scratches)~\cite{chen2025obb,zhao2024small}.
\item \textbf{NMS Suppression of Clustered Flaws:} When multiple defects occur in close spatial proximity along container edges, horizontal bounding boxes overlap excessively. During standard Non-Maximum Suppression (NMS), overlapping boxes suppress valid adjacent defect detections, generating false-negative escapes~\cite{yang2024gradient}.
\end{itemize}

%------------------------------------------------------------------
\section{Aim and Motivation}
\label{sec:aim}

\subsection{Aim}
\label{subsec:aim}

To develop, optimize, and deploy a vision-based defect detection system for industrial quality inspection using an oriented bounding box deep learning model (YOLOv8-OBB) for high-speed, sub-millimeter surface defect detection on thermoformed industrial packaging casings.

\subsection{Background}
\label{subsec:aimbackground}

Oriented Bounding Box (OBB) representation extends conventional object detection by predicting an orientation angle parameter $\theta$ in addition to centroid coordinates and dimensions $[x, y, w, h, \theta]$. While widely explored in aerial satellite imagery, OBB architectures have rarely been adapted for high-speed, resource-constrained industrial manufacturing cells where inference latency must remain strictly below 25~ms~\cite{wang2024rgs,huang2026msfe}. By tightly enclosing flaws along their principal axes of inertia, OBB eliminates extraneous background pixels, preserves high-frequency edge gradients, and provides exact geometric coordinates for downstream robotic pick-and-place actuators~\cite{zhang2024robotic}.

\subsection{Motivation for the Project}
\label{subsec:motivation}

The motivation for this project stems from urgent industrial and economic demands. Contaminated or compromised food packaging triggers catastrophic commercial repercussions: nationwide product recalls costing millions of dollars, severe regulatory penalties, and reputational damage. A single undetected 0.2~mm pinhole can spoil an entire production batch within 48 hours of hermetic sealing. By deploying an intelligent, deterministic edge vision inspection cell operating at over 40 frames per second with near-zero defect escape (less than 0.3\%), manufacturing enterprises can eliminate human inspection labor overheads, reduce scrap rates, and ensure complete packaging integrity~\cite{meng2024automated,sun2026edgedeployable}.

\subsection{Expected Contributions of the Project}
\label{subsec:contributions}

This project delivers four key technical and architectural contributions to the field of automated optical quality engineering:

\begin{enumerate}
\item \textbf{Orientation-Aware Deep Architecture:} Formulation and implementation of an optimized YOLOv8-OBB deep neural network incorporating decoupled oriented bounding box heads, Gaussian Wasserstein Distance (GWD) regression loss, and Varifocal classification loss tailored for high-aspect-ratio packaging defects.
\item \textbf{Micro-Defect Augmentation Theory:} Discovery and rigorous empirical validation of micro-defect augmentation suppression: demonstrating that eliminating spatial mosaic transformations and restricting geometric translation/scaling prevents the destruction of slender sub-pixel flaw features, elevating micro-scratch detection mAP50 by an astonishing 83.24 percentage points (from 12.61\% to 95.85\%).
\item \textbf{Comprehensive Industrial Flaw Benchmark:} Compilation of an industrial packaging benchmark dataset comprising 5,065 high-resolution images annotated across nine defect categories with ground-truth rotated bounding polygons, establishing an open benchmark for packaging quality inspection.
\item \textbf{Real-Time Robotic Edge Integration:} Development of an integrated edge deployment pipeline utilizing ONNX-Slim optimization, TensorRT acceleration, and a ROS2-PLC real-time middleware synchronization architecture that achieves an end-to-end cycle latency of 24.0~ms (41.6~FPS) on edge industrial hardware.
\end{enumerate}

%------------------------------------------------------------------
\section{General Problem Statement}
\label{sec:problem}

In high-speed industrial packaging manufacturing lines operating at conveyance velocities of 0.8 to 1.2~m/s, conventional human inspection and horizontal automated optical inspection (AOI) frameworks fail to reliably detect, classify, and localize slender, diagonally oriented, and sub-millimeter surface flaws (such as micro-scratches, pinholes, and shrivel wrinkles) due to extreme background feature dilution, illumination sensitivity, and severe execution latency bottlenecks. This leads to prohibitive defect escape rates exceeding 14\%, frequent food packaging barrier failure, and massive commercial recall risks. There is an urgent engineering need for a deterministic, real-time, vision-based defect detection system for industrial quality inspection that achieves greater than 95\% mAP50 across all defect classes within a strict 25~ms cycle time.

%------------------------------------------------------------------
\section{Research Question}
\label{sec:rq}

How can deep oriented bounding box regression architectures and hyperparameter augmentation strategies be optimized and integrated with embedded edge accelerators to achieve real-time, sub-millimeter surface flaw localization on high-speed industrial packaging conveyors without suffering from small-object feature dilution or exceeding factory latency budgets?

\subsection{Research Objectives}
\label{subsec:objectives}

To systematically address the research question, five measurable engineering and scientific objectives were defined and executed:

\begin{enumerate}
\item Curate, clean, and annotate an exhaustive industrial dataset of 5,065 high-resolution packaging casing images across nine flaw classes using rotated bounding boxes $[x, y, w, h, \theta]$ to establish ground-truth orientation baselines.
\item Design and construct an industrial optomechanical vision enclosure featuring coaxial diffuse ring illumination, high-speed optical triggering, and telecentric optics to eliminate specular reflections and optical distortion on curved polymer casings.
\item Formulate, adapt, and train an optimized YOLOv8-OBB model incorporating decoupled orientation heads, Gaussian Wasserstein Distance loss, and Varifocal classification loss to resolve boundary discontinuity and feature dilution.
\item Conduct systematic ablation studies to investigate the impact of spatial mosaic augmentation, scale transformations, and adaptive optimizers (AdamW vs. SGD-M) on micro-defect detection fidelity.
\item Serialize, optimize, and deploy the neural network on an edge computing node via ONNX/TensorRT, establishing real-time ROS2 middleware synchronization with downstream pneumatic ejection actuators to achieve continuous sorting under 25~ms latency.
\end{enumerate}

%------------------------------------------------------------------
\section{Project Report Organization}
\label{sec:org}

The remainder of this Major Project Report is structured into four cohesive chapters:

\begin{itemize}
\item \textbf{Chapter 2 (Literature Survey on Vision-Based Defect Detection System for Industrial Quality Inspection):} Presents a comprehensive systematic literature survey conducted in accordance with PRISMA guidelines. Evaluates 30 landmark peer-reviewed studies across deep learning AOI, oriented bounding box models, small object detection, and edge robotics, culminating in a detailed literature review table and explicit research gap analysis.
\item \textbf{Chapter 3 (Methodology on Vision-Based Defect Detection System for Industrial Quality Inspection):} Details the end-to-end system methodology, optomechanical enclosure design, 5,065-image dataset taxonomy, YOLOv8-OBB network architecture, mathematical loss derivations (GWD, VFL, DFL), rotated polygon NMS clipping, micro-defect augmentation suppression theory, and ROS2/PLC edge integration.
\item \textbf{Chapter 4 (Results and Discussion):} Provides extensive empirical findings, quantitative per-class performance comparisons, stepwise ablation analysis (M1--M5), training convergence curves, normalized confusion matrix, visual qualitative detections, edge latency profiling across hardware targets, cross-architecture benchmarks, economic ROI analysis, and failure mode mitigation protocols.
\item \textbf{Chapter 5 (Conclusion \& Future Scope):} Synthesizes the overall research achievements against each stated objective, lists the key quantitative project outcomes, outlines future directions for short-wave infrared (SWIR) imaging and INT8 edge quantization, and documents research publications.
\end{itemize}
